\subsection{Sensorem}

Sensorem er en svensk virksomhed, der har specialiseret sig i tryghedsalarmer til ældre. Baggrunden for virksomheden er den udfordring, som både Sverige og Danmark står overfor: en voksende ældre befolkning. Sensorems mission er at skabe løsninger, der giver ældre større tryghed i hverdagen, så de kan blive længst muligt i deres eget hjem og i mange tilfælde helt undgå behovet for plejehjem.

Virksomheden har udviklet et diskret ur, hvor alarmen er integreret, så den ikke er synlig. Uret har en knap, som brugeren kan trykke på, hvis der opstår behov for hjælp. Derudover er det udstyret med en automatisk faldalarm, der registrerer, hvis brugeren falder, og straks aktiverer alarmen. I sådanne tilfælde bliver de pårørende ringet op via urets indbyggede højtalertelefon.

Urets højtaler gør det muligt for de pårørende at kommunikere direkte med brugeren via en app. Hvis ingen af de pårørende svarer, tager Sensorems alarmcentral over og håndterer situationen.

Uret er koblet op på mobilnetværket, hvilket betyder, at det fungerer overalt, også uden for hjemmet. Skulle der ske en hændelse på uden for brugerens bopæl, sendes signalet stadig til de pårørende. Desuden er uret udstyret med GPS, så de pårørende altid kan se brugerens position i realtid. 

En vigtig fordel er, at tryghedsalarmen ikke kræver tilknytning til hjemmeplejen. Alarmen går direkte til de pårørende. Hvis de pårørende ikke responderer på alarmen, sendes den automatisk videre til Sensorems alarmcentral, der kan tilkalde en ambulance.

Automatisk svar på opkald

Vandtæt

Inaktivitetsalarm 

Medicinpåmindelse